本实验借助 MATLAB 和 Simulink 研究了基本传递环节（P、I、PT1、PT2）对阶跃信号和正弦信号的动态行为。
仿真清楚展示了各系统的特有性质：P 环节无延迟响应，I 环节呈积分上升，而 PT1 和 PT2 环节表现出延迟特性。
对于 PT2 环节，固有频率和阻尼系数主要决定响应速度以及超调倾向。
阻尼越小，阶跃响应中的振荡和超调越明显，Bode 图中的共振峰也越突出；阻尼增大后，系统响应更加平滑但可能变慢。
在阻尼保持不变时，固有频率越大，极点离原点越远，系统响应越快，频率响应的特征点也向高频方向移动。

通过 Bode 图和零极点图进行的补充分析直观说明了频率响应和极点几何位置如何决定时域行为与系统稳定性。
由此可以学习如何在实践中解释线性系统的理论模型。对反应时间和振荡倾向如何由参数系统性预测与调节的认识，为复杂控制回路的设计提供了重要且直观的基础。
